\documentclass[aspectratio=169,xcolor=dvipsnames]{beamer}
\usetheme{Madrid}\usecolortheme{seahorse}
\usepackage{amsmath,amssymb,listings,tikz,booktabs,hyperref}
\usetikzlibrary{arrows.meta,positioning,shapes}
\lstset{language=Python,basicstyle=\small\ttfamily,keywordstyle=\color{blue}\bfseries,
  commentstyle=\color{OliveGreen},stringstyle=\color{red},frame=single,breaklines=true}

\title[Algorithms -- Week 7]{\textbf{Thiết kế và Phân tích Thuật toán}\\
  \large Week 7: Dynamic Programming I}
\subtitle{Introduction, Knapsack, LCS, Memoisation \\ TOAE209/TOAH209 -- Foreign Trade University}
\author{Department of Technology \& Data Science}
\date{Semester 2024--2025}
\AtBeginSection[]{\begin{frame}<beamer>{Outline}\tableofcontents[currentsection]\end{frame}}

\begin{document}
\begin{frame}\titlepage\end{frame}
\begin{frame}{Outline}\tableofcontents\end{frame}

\section{Motivation}
\begin{frame}{Why Dynamic Programming?}
  \begin{columns}
    \column{0.55\textwidth}
    \begin{alertblock}{The Problem with Naive Recursion}
      Fibonacci: $T(n) = T(n-1)+T(n-2)+1$\\
      $\Rightarrow T(n) = \Theta(2^n)$\\[4pt]
      $\text{fib}(5) = \text{fib}(4)+\text{fib}(3)$\\
      $\text{fib}(4) = \text{fib}(3)+\text{fib}(2)$\\
      $\text{fib}(3)$ computed \textbf{twice!}
    \end{alertblock}
    \column{0.4\textwidth}
    \begin{exampleblock}{DP Solution}
      Store computed results.\\
      Each subproblem solved \textbf{once}.\\
      $\Theta(n)$ time, $\Theta(n)$ space.
    \end{exampleblock}
  \end{columns}
  \vspace{0.5em}
  \begin{block}{Two Key Properties for DP}
    \begin{enumerate}
      \item \textbf{Overlapping subproblems}: same subproblems recur many times
      \item \textbf{Optimal substructure}: optimal solution built from optimal sub-solutions
    \end{enumerate}
  \end{block}
\end{frame}

\section{Historical Context}
\begin{frame}{History of Dynamic Programming}
  \begin{description}
    \item[1950s] Richard Bellman coins ``dynamic programming'' (intentionally vague name!)
    \item[1953] Bellman's principle of optimality
    \item[1957] Floyd-Warshall all-pairs shortest paths
    \item[1965] Needleman-Wunsch sequence alignment (bioinformatics)
    \item[1970s] DP used in compiler design (CYK parsing)
    \item[1990s] DP in speech recognition (HMMs), finance (options pricing)
    \item[2010s] Deep Reinforcement Learning = DP on massive state spaces
  \end{description}
\end{frame}

\section{Memoisation vs Tabulation}
\begin{frame}[fragile]{Top-Down: Memoisation}
  \begin{lstlisting}
# Memoised Fibonacci: O(n) time, O(n) space
def fib(n, memo={}):
    if n in memo: return memo[n]     # Cache hit
    if n <= 1: return n
    memo[n] = fib(n-1, memo) + fib(n-2, memo)
    return memo[n]
  \end{lstlisting}
  \begin{block}{Key Points}
    \begin{itemize}
      \item Recursive structure preserved (natural to write)
      \item Only computes subproblems actually needed
      \item Risk: stack overflow for large $n$
    \end{itemize}
  \end{block}
\end{frame}

\begin{frame}[fragile]{Bottom-Up: Tabulation}
  \begin{lstlisting}
# Tabulation Fibonacci: O(n) time, O(1) space
def fib_tab(n):
    if n <= 1: return n
    a, b = 0, 1
    for _ in range(2, n+1):
        a, b = b, a + b
    return b
  \end{lstlisting}
  \begin{block}{Key Points}
    \begin{itemize}
      \item Fill table in bottom-up order (base cases first)
      \item Often more efficient (no function call overhead)
      \item Can optimise space by only keeping needed rows
    \end{itemize}
  \end{block}
\end{frame}

\section{Classic DP Problems}
\begin{frame}[fragile]{0/1 Knapsack Problem}
  \textbf{Problem:} Given $n$ items with weights $w_i$ and values $v_i$, and capacity $W$,
  find the maximum value subset with total weight $\le W$.\\[6pt]
  \textbf{Subproblem:} $dp[i][w]$ = max value using items $1..i$ with capacity $w$.\\[6pt]
  \textbf{Recurrence:}
  $$dp[i][w] = \begin{cases}
    dp[i-1][w] & \text{if } w_i > w \text{ (can't include)} \\
    \max(dp[i-1][w],\; dp[i-1][w-w_i]+v_i) & \text{otherwise}
  \end{cases}$$
  \textbf{Complexity:} $\Theta(nW)$ time, $\Theta(nW)$ space
  \begin{alertblock}{Note}
    This is \textbf{pseudo-polynomial}: depends on the value $W$, not $\log W$.
    The 0/1 knapsack is NP-hard in general!
  \end{alertblock}
\end{frame}

\begin{frame}[fragile]{Longest Common Subsequence (LCS)}
  \textbf{Recurrence:}
  $$dp[i][j] = \begin{cases}
    0 & \text{if } i=0 \text{ or } j=0 \\
    dp[i-1][j-1]+1 & \text{if } s[i] = t[j] \\
    \max(dp[i-1][j], dp[i][j-1]) & \text{otherwise}
  \end{cases}$$
  \textbf{Example:} LCS(``ABCBDAB'', ``BDCAB'') = ``BCAB'' (length 4)\\[6pt]
  \begin{block}{LCS Table Fragment}
    \begin{tabular}{c|ccccc}
      & B & D & C & A & B \\ \hline
      A & 0 & 0 & 0 & 1 & 1 \\
      B & 1 & 1 & 1 & 1 & 2 \\
      C & 1 & 1 & 2 & 2 & 2 \\
    \end{tabular}
  \end{block}
\end{frame}

\section{Further Reading}
\begin{frame}{Further Reading \& Advanced Topics}
  \begin{block}{References}
    \begin{itemize}
      \item CLRS Ch. 15: Dynamic Programming (full treatment)
      \item MIT 6.006 Lecture Notes on DP
    \end{itemize}
  \end{block}
  \begin{exampleblock}{Applications}
    \begin{itemize}
      \item Bioinformatics: sequence alignment (Smith-Waterman)
      \item Natural Language Processing: CYK parsing, Viterbi algorithm
      \item Computer Vision: Viterbi HMM, CRF inference
      \item Finance: Black-Scholes option pricing via DP
      \item Reinforcement Learning: value iteration, Q-learning
    \end{itemize}
  \end{exampleblock}
\end{frame}

\begin{frame}{Week 7 Summary}
  Review the lecture notes, complete \texttt{starter.py}, and read the assigned chapters.\\[6pt]
  \begin{block}{Self-Check Questions}
    \begin{enumerate}
      \item Can you implement the key algorithms from scratch?
      \item Can you prove correctness using a loop invariant?
      \item Can you derive and solve the time complexity recurrence?
      \item Can you identify the best and worst case inputs?
    \end{enumerate}
  \end{block}
\end{frame}
\end{document}
